%% Journal of Scientific Computing submission -- frame for main.tex.
%% main.tex is GENERATED by flatten.py from this frame + sections/*.tex (verbatim copies of
%% paper2d/sections). Edit this frame or the section files, then rerun flatten.py.
%% Springer Nature LaTeX template (sn-jnl.cls, "Version 2.1 April 2023" per the template's
%% sample file), obtained from the rstudio/rticles GitHub mirror because the journal's own
%% download (media.springer.com) is not reachable from this environment -- see GUIDELINES.txt.
\PassOptionsToPackage{hyphens}{url}
%% amsthm must be loaded before sn-jnl.cls so that the class defines its theorem styles
%% (thmstyleone, thmstylethree) and its proof environment (it tests \@ifpackageloaded{amsthm}).
\RequirePackage{amsmath}\RequirePackage{amsthm}
\documentclass[pdflatex,sn-mathphys-num]{sn-jnl}

\usepackage{graphicx}%
\usepackage{amsmath,amssymb,amsfonts}%
\usepackage{amsthm}%
\usepackage{mathtools,booktabs,array,enumitem}%
\usepackage{xcolor}%
\usepackage{url}
\usepackage[expansion=false]{microtype}
\usepackage{tikz}
\usetikzlibrary{calc}

\setlist{leftmargin=2em}
\emergencystretch=2em
\raggedbottom

%% Theorem environments: same names, same shared counter numbered within sections as in
%% paper2d (Theorem 5.3 etc.), in the template's styles.
\theoremstyle{thmstyleone}
\newtheorem{theorem}{Theorem}[section]
\newtheorem{proposition}[theorem]{Proposition}
\newtheorem{lemma}[theorem]{Lemma}
\newtheorem{corollary}[theorem]{Corollary}
\theoremstyle{thmstylethree}
\newtheorem{definition}[theorem]{Definition}
\newtheorem{remark}[theorem]{Remark}
\newtheorem{assumption}[theorem]{Assumption}

%%MACROS%%
\begin{document}

\title[Scott--Vogelius--Nitsche on polygonal boundaries]{Exactly divergence-free Stokes elements with Nitsche's method on a polygonally approximated boundary: the missing pressure traction, Scott's rate, and the zero-gradient lock}

\author*[1]{\fnm{Jaideep Sai} \sur{Padhi}}\email{jpadhi@purdue.edu}
%% FILL: ORCID if available, e.g. \author*[1]{\fnm{Jaideep Sai} \sur{Padhi}~\orcidlink{0000-0000-0000-0000}}
%%       (and \usepackage{orcidlink}); also enter it in Editorial Manager.
%% TBD: C. de S. Cavalcante may join as a co-author (per the joint AML letter); not added here.

\affil*[1]{\orgdiv{Department of Computer Science}, \orgname{Purdue University}, \orgaddress{\city{West Lafayette}, \postcode{47907}, \state{IN}, \country{USA}}}

%% Abstract shortened to the journal's 150--250 word limit (the paper2d abstract has ~410 words);
%% no claim is added, every claim below is stated in the paper.
\abstract{Scott--Vogelius elements compute pointwise divergence-free velocities, but when a curved no-slip wall is replaced by an inscribed polygon discrete velocities can be forced to have zero gradient at polygon vertices, and with strongly imposed no-slip the $H^1$ error can degrade to order $h^{1/2}$. Gjerde and Scott imposed the wall condition weakly by Nitsche's method on the polygon, and Scott asked for a proof that the error behaves like $h_\Gamma^{3/2}+h_\Omega^k$ and, if not, why not. We answer the question in two dimensions. For the method as printed, the estimate holds if and only if the pressure is constant along the wall: the Nitsche form omits the pressure traction $-pn$, and the velocity leaks through the wall with normal velocity $(h/\mu)(p-\bar p)$. We prove sharp rates for this error and identify its leading term. Restoring the traction with its wall mean removed, a device of Frachon, Nilsson and Zahedi, gives an exactly divergence-free method; we prove that it attains $h_\Gamma^{3/2}+h_\Omega^k$ for every $k\ge4$, uniformly in the penalty for flux-corrected data, and that $h_\Gamma^{3/2}$ is sharp whenever the wall shear does not vanish identically and $h_\Gamma$ is small relative to the smallest wall-triangle angle. The penalty must scale with the ratio of bulk to boundary mesh size (sufficient on every mesh, necessary on broad classes of meshes). Under the Guzm\'an--Scott angle condition strong imposition attains the same rate; the zero-gradient lock is caused by wall vertices lying in only two triangles. Computations confirm the rates and constants.}

\keywords{Stokes equations, Scott--Vogelius elements, Nitsche's method, curved boundaries, divergence-free finite elements, penalty parameter}

\pacs[MSC Classification]{65N30 (primary), 65N12, 65N15, 76D07}

\maketitle

%%BODY%%
\backmatter

%%ACK%%
%%DECLARATIONS%%
\bibliography{refs}

\end{document}
